\documentclass{article}
\usepackage[T2A]{fontenc}
\usepackage[utf8]{inputenc}
\usepackage{amsthm}
\usepackage{amsmath}
\usepackage{amssymb}
\usepackage{amsfonts}
\usepackage{mathrsfs}
\usepackage[12pt]{extsizes}
\usepackage{fancyvrb}
\usepackage{indentfirst}
\usepackage[
left=2cm, right=2cm, top=2cm, bottom=2cm, headsep=0.2cm, footskip=0.6cm, bindingoffset=0cm
]{geometry}
\usepackage[english,russian]{babel}


\begin{document}
	\section*{Вариант 7}
	Если требуется улучшить качество переходных процессов при некотором фиксированном  
	\( s = s_0 \), то в этом случае целесообразно выполнить параметрический синтез, т. е. выбор значений параметров обратных связей, на основе минимизации функции \( f(p, s_0) \), где  
	
	\[
	f(p, s) = \left( \|R_A(0, p_0, s)\|^{-2} + \|R_A(0, p, s)\|^{-2} \right) \int_{0}^{\infty} f(p, s, \omega) \, d\omega,
	\]
	
	\[
	f(p, s, \omega) = \|R_A(\omega, p, s) - R_A(0, p, s)R_A^*(\omega)\|^2 
	+ c_1 \|R'_A(\omega, p, s) - R_A(0, p, s)R_A^{*'}(\omega)\|^2 
	+ c_2 \|R''_A(\omega, p, s) - R_A(0, p, s)R_A^{*''}(\omega)\|^2, \quad (') = d()/d\omega,
	\]
	
	\[
	R_{A\nu j}(\omega, p, s) = 
	\begin{cases} 
		\Re \Phi_{\nu j}(i\omega, p, s), & A_j(p, s) \neq 0, \\ 
		\sqrt{1 + \omega^2} \;\Re[\Psi_{\nu j}(i\omega, p, s)], & A_j(p, s) = 0,
	\end{cases}
	\]
\end{document}

